\documentclass[11pt, a4paper]{article}

% ── Packages ─────────────────────────────────────────────────────────────────
\usepackage[utf8]{inputenc}
\usepackage[T1]{fontenc}
\usepackage[spanish]{babel}
\usepackage{geometry}
\usepackage{microtype}
\usepackage{booktabs}
\usepackage{longtable}
\usepackage{array}
\usepackage{xcolor}
\usepackage{colortbl}
\usepackage{amsmath}
\usepackage{amssymb}
\usepackage{hyperref}
\usepackage{fancyhdr}
\usepackage{titlesec}
\usepackage{enumitem}
\usepackage{parskip}
\usepackage{mdframed}
\usepackage{tikz}
\usepackage{graphicx}
\usetikzlibrary{arrows.meta, positioning, calc, backgrounds, decorations.pathmorphing}

% ── Geometría ────────────────────────────────────────────────────────────────
\geometry{top=2.5cm, bottom=2.5cm, left=2.8cm, right=2.8cm, headheight=14pt}

% ── Colores ──────────────────────────────────────────────────────────────────
\definecolor{tlvblue}{RGB}{31, 56, 100}
\definecolor{tlvmid}{RGB}{46, 80, 144}
\definecolor{tlvlight}{RGB}{214, 228, 248}
\definecolor{tlvgray}{RGB}{240, 244, 250}
\definecolor{tlvrule}{RGB}{200, 200, 200}
\definecolor{tlvred}{RGB}{180, 40, 40}
\definecolor{tlvgreen}{RGB}{30, 120, 60}
\definecolor{tlvorange}{RGB}{200, 100, 0}
\definecolor{steelgray}{RGB}{120, 130, 145}
\definecolor{zonered}{RGB}{255, 220, 220}
\definecolor{zoneyellow}{RGB}{255, 243, 200}
\definecolor{zonegreen}{RGB}{220, 245, 220}

% ── Hyperref ─────────────────────────────────────────────────────────────────
\hypersetup{colorlinks=true, linkcolor=tlvmid, urlcolor=tlvmid,
  pdftitle={TLV-1 DDR-003}, pdfauthor={Programa TLV}}

% ── Encabezados ──────────────────────────────────────────────────────────────
\pagestyle{fancy}
\fancyhf{}
\fancyhead[L]{\small\color{tlvblue}\textbf{TLV-1} \textbar\ DDR-003}
\fancyhead[R]{\small\color{tlvblue} Plataforma de Lanzamiento y Erector Hidr\'{a}ulico}
\fancyfoot[C]{\small\color{gray} P\'{a}gina \thepage}
\renewcommand{\headrulewidth}{0.5pt}
\renewcommand{\headrule}{\hbox to\headwidth{\color{tlvblue}\leaders\hrule height \headrulewidth\hfill}}

% ── Títulos ───────────────────────────────────────────────────────────────────
\titleformat{\section}{\large\bfseries\color{tlvblue}}{\thesection.}{0.6em}{}
  [\vspace{2pt}\color{tlvrule}\titlerule]
\titleformat{\subsection}{\normalsize\bfseries\color{tlvmid}}{\thesubsection.}{0.5em}{}
\titlespacing*{\section}{0pt}{18pt}{8pt}
\titlespacing*{\subsection}{0pt}{12pt}{4pt}

% ── Columnas ─────────────────────────────────────────────────────────────────
\newcolumntype{L}[1]{>{\raggedright\arraybackslash}p{#1}}
\newcolumntype{C}[1]{>{\centering\arraybackslash}p{#1}}

% ── Comandos ─────────────────────────────────────────────────────────────────
\newcommand{\status}[1]{\colorbox{tlvgreen}{\color{white}\small\bfseries\strut\ #1\ }}
\newcommand{\cmark}{\textcolor{tlvgreen}{\textbf{\checkmark}}}
\newcommand{\xmark}{\textcolor{tlvred}{\textbf{$\times$}}}

% ─────────────────────────────────────────────────────────────────────────────
\begin{document}

% ── PORTADA ──────────────────────────────────────────────────────────────────
\begin{titlepage}
  \centering\vspace*{3cm}
  {\large\color{gray}\bfseries PROGRAMA TLV}\\[0.4cm]
  {\Huge\color{tlvblue}\bfseries Test Launch Vehicle 1}\\[0.3cm]
  {\Large\color{tlvmid} DDR-003 --- Design Decision Record}\\[2cm]
  \color{tlvrule}\hrule height 0.8pt\vspace{0.6cm}
  \begin{tabular}{L{4cm} L{9cm}}
    \rowcolor{tlvgray}
    \textbf{Documento}  & DDR-003 \\
    \textbf{T\'{i}tulo} & Plataforma de Lanzamiento, Erector Hidr\'{a}ulico y Zona NOTAM \\
    \rowcolor{tlvgray}
    \textbf{Subsistema} & Infraestructura de Tierra / Operaciones de Lanzamiento \\
    \textbf{Estado}     & \status{APROBADO} \\
    \rowcolor{tlvgray}
    \textbf{Revisi\'{o}n} & 1.0 \\
    \textbf{Referencias} & DDR-001 (Separaci\'{o}n Nariz), DDR-002 (Edificio AIT) \\
  \end{tabular}
  \vspace{0.6cm}\color{tlvrule}\hrule height 0.8pt
  \vfill
  {\small\color{gray} Programa TLV --- Documento de Decisiones de Dise\~{n}o}
\end{titlepage}

% ── 1. CONTEXTO ──────────────────────────────────────────────────────────────
\section{Contexto}

La plataforma de lanzamiento es la infraestructura f\'{i}sica desde la cual el TLV-1 es erigido y lanzado. Este DDR documenta tres decisiones interrelacionadas: la configuraci\'{o}n y dimensiones de la plataforma, el dise\~{n}o del erector hidr\'{a}ulico de pivote, y el radio de la zona de exclusi\'{o}n NOTAM derivado del an\'{a}lisis cin\'{e}tico de riesgo.

Los par\'{a}metros de vuelo utilizados en este documento provienen de la simulaci\'{o}n num\'{e}rica del perfil de vuelo del TLV-1, integrada con paso $\Delta t = 0.1\ \text{s}$ incluyendo arrastre aerodin\'{a}mico variable con altitud.

% ── 2. PARÁMETROS DE VUELO RELEVANTES ────────────────────────────────────────
\section{Par\'{a}metros de Vuelo Relevantes}

\begin{longtable}{L{5cm} C{3.5cm} L{5cm}}
  \rowcolor{tlvblue}
  \color{white}\textbf{Par\'{a}metro} &
  \color{white}\textbf{Valor} &
  \color{white}\textbf{Fuente} \\
  \endfirsthead
  \rowcolor{tlvblue}
  \color{white}\textbf{Par\'{a}metro} &
  \color{white}\textbf{Valor} &
  \color{white}\textbf{Fuente} \\
  \endhead
  \rowcolor{tlvgray}
  Masa total al despegue ($m_0$)    & 369 kg          & Dise\~{n}o TLV-1 \\
  Masa seca ($m_{dry}$)             & 250 kg          & Estimado (pendiente CAD) \\
  \rowcolor{tlvgray}
  Empuje                            & 7{,}310 N       & Calculado \\
  TWR inicial                       & 2.02            & Calculado \\
  \rowcolor{tlvgray}
  Aceleraci\'{o}n inicial           & 10 m/s$^2$      & Requerimiento aviónica \\
  Tiempo de quemado ($t_{burn}$)    & 33.5 s          & Calculado \\
  \rowcolor{tlvgray}
  Velocidad en burnout              & 364.4 m/s       & Simulaci\'{o}n num\'{e}rica \\
  Altitud en burnout                & 6{,}248 m       & Simulaci\'{o}n num\'{e}rica \\
  \rowcolor{tlvgray}
  Tiempo a Mach 1                   & 29.2 s          & Simulaci\'{o}n num\'{e}rica \\
  Altitud en Mach 1                 & 4{,}770 m       & Simulaci\'{o}n num\'{e}rica \\
  \rowcolor{tlvgray}
  Apogeo                           & 11{,}450 m      & Simulaci\'{o}n num\'{e}rica \\
  Energ\'{i}a cin\'{e}tica m\'{a}x.\ (burnout) & $\approx$16.6 MJ & Simulaci\'{o}n num\'{e}rica \\
\end{longtable}

% ── 3. PLATAFORMA DE LANZAMIENTO ─────────────────────────────────────────────
\section{Plataforma de Lanzamiento}

\subsection{Filosofía de dise\~{n}o}

La plataforma debe ser lo m\'{a}s ligera y compacta posible, coherente con la escala del veh\'{i}culo. No es necesaria una estructura masiva: el TLV-1 es un veh\'{i}culo de 369 kg con un empuje de 7{,}310 N y un tiempo de exposici\'{o}n al chorro muy breve antes del despegue.

El chorro de la tobera impacta la plataforma durante aproximadamente 0.5 s antes de que el cohete se eleve suficientemente. Con una placa sacrificial reemplazable en la zona de impacto, la estructura principal no sufre da\~{n}o t\'{e}rmico significativo.

\subsection{Especificaciones}

\begin{longtable}{L{5cm} C{3.5cm} L{5cm}}
  \rowcolor{tlvblue}
  \color{white}\textbf{Par\'{a}metro} &
  \color{white}\textbf{Valor} &
  \color{white}\textbf{Justificaci\'{o}n} \\
  \endfirsthead
  \rowcolor{tlvblue}
  \color{white}\textbf{Par\'{a}metro} &
  \color{white}\textbf{Valor} &
  \color{white}\textbf{Justificaci\'{o}n} \\
  \endhead
  \rowcolor{tlvgray}
  Dimensiones en planta      & $2.5\ \times\ 2.5$ m   & Huella m\'{i}nima para estabilidad del erector \\
  Altura sobre el suelo      & 0.5--0.8 m             & Espacio para tobera y umbilicales \\
  Material estructura        & Acero estructural S275  & Econom\'{i}a, soldabilidad, disponibilidad \\
  \rowcolor{tlvgray}
  Placa sacrificial          & Acero 8 mm, 0.5$\times$0.5 m & Zona de impacto del chorro; reemplazable \\
  Carga de dise\~{n}o        & $\geq 5{,}000$ N        & Empuje + factor din\'{a}mico de erector \\
  \rowcolor{tlvgray}
  Anclaje al suelo           & 4 pernos de anclaje M24 & Retenci\'{o}n del cohete en cuenta regresiva \\
  Conexiones umbilicales     & 2 puntos laterales      & Alimentaci\'{o}n el\'{e}ctrica + telemetr\'{i}a de tierra \\
  \rowcolor{tlvgray}
  Deflector de gases         & No requerido            & $\Delta T \approx 25°C$ en placa; v\'{e}ase \S4 \\
\end{longtable}

\subsection{Verificaci\'{o}n t\'{e}rmica: placa de impacto}

El flujo de calor por convecci\'{o}n en la zona de impacto directo del chorro es:

\begin{equation}
  q'' \approx 2\ \text{MW/m}^2
\end{equation}

El aumento de temperatura en una placa de acero de 8 mm durante el tiempo de exposici\'{o}n estimado ($t_{exp} \approx 0.5\ \text{s}$):

\begin{equation}
  \Delta T = \frac{q'' \cdot t_{exp}}{\rho \cdot c_p \cdot e}
  = \frac{2 \times 10^6 \times 0.5}{7850 \times 500 \times 0.008}
  = \frac{10^6}{31{,}400} \approx \mathbf{32°C}
\end{equation}

La temperatura final de la placa ($\approx 52°C$ desde temperatura ambiente) est\'{a} muy por debajo del l\'{i}mite del acero estructural. El deflector queda descartado para TLV-1.

\subsection{Esquema de la plataforma}

\begin{center}
\begin{tikzpicture}[scale=1.0]

  % Base de la plataforma
  \fill[steelgray!30] (-3.5, 0) rectangle (3.5, 0.6);
  \draw[steelgray, thick] (-3.5, 0) rectangle (3.5, 0.6);
  \node[font=\small, color=steelgray] at (0, 0.3) {Estructura acero S275};

  % Patas
  \foreach \x in {-3.0, 3.0} {
    \fill[steelgray!50] (\x-0.15, -0.8) rectangle (\x+0.15, 0);
    \draw[steelgray] (\x-0.15, -0.8) rectangle (\x+0.15, 0);
  }

  % Pernos de anclaje
  \foreach \x in {-3.0, 3.0} {
    \fill[tlvblue] (\x-0.08, -1.1) rectangle (\x+0.08, -0.8);
    \draw[tlvblue, thick] (\x, -1.1) -- (\x, -1.4);
  }
  \node[font=\tiny, color=tlvblue] at (0, -1.25) {Pernos de anclaje M24};

  % Placa sacrificial
  \fill[tlvred!40] (-0.8, 0.6) rectangle (0.8, 0.72);
  \draw[tlvred, thick] (-0.8, 0.6) rectangle (0.8, 0.72);
  \node[font=\tiny, color=tlvred] at (2.2, 0.66) {Placa sacrificial 8 mm};
  \draw[tlvred, thin, -latex] (1.5, 0.66) -- (0.85, 0.66);

  % Pivote del erector
  \fill[tlvblue] (-3.2, 0.6) circle (0.15);
  \draw[tlvblue, thick] (-3.2, 0.6) circle (0.15);
  \node[font=\tiny, color=tlvblue, align=center] at (-3.2, 1.1) {Pivote\\erector};

  % Cohete (vertical, simplificado)
  \fill[tlvmid!30] (-0.22, 0.72) rectangle (0.22, 5.72);
  \draw[tlvmid, thick] (-0.22, 0.72) rectangle (0.22, 5.72);
  % Nariz
  \fill[tlvblue]
    (-0.22, 5.72) -- (0, 6.22) -- (0.22, 5.72) -- cycle;
  % Aletas esquem.
  \fill[steelgray!60]
    (-0.22, 0.72) -- (-0.7, 0.3) -- (-0.7, 0.72) -- cycle;
  \fill[steelgray!60]
    (0.22, 0.72) -- (0.7, 0.3) -- (0.7, 0.72) -- cycle;
  \node[font=\tiny, color=tlvmid] at (1.2, 3.2) {TLV-1 (5 m)};

  % Chorro de la tobera
  \fill[tlvred!20]
    (-0.15, 0.72) -- (-0.5, -0.1) -- (0, 0.1) -- (0.5, -0.1) -- (0.15, 0.72) -- cycle;
  \draw[tlvred!60, thick, decorate, decoration={snake, amplitude=1pt, segment length=4pt}]
    (0, 0.72) -- (0, -0.05);
  \node[font=\tiny, color=tlvred] at (1.3, 0.2) {Chorro APCP};

  % Umbilicales
  \draw[tlvorange, thick, dashed] (0.22, 2.0) -- (1.5, 2.0) -- (1.5, -0.5);
  \draw[tlvorange, thick, dashed] (-0.22, 1.5) -- (-1.5, 1.5) -- (-1.5, -0.5);
  \node[font=\tiny, color=tlvorange] at (2.3, 1.5) {Umbilicales};

  % Cotas
  \draw[latex-latex, gray] (0.35, 0.72) -- (0.35, 5.72);
  \node[font=\tiny, color=gray, rotate=90] at (0.65, 3.2) {5000 mm};

  \draw[latex-latex, gray] (-3.5, -1.5) -- (3.5, -1.5);
  \node[font=\tiny, color=gray] at (0, -1.75) {2500 mm};

\end{tikzpicture}
\end{center}

% ── 4. ERECTOR HIDRÁULICO ─────────────────────────────────────────────────────
\section{Erector Hidr\'{a}ulico de Pivote}

\subsection{Concepto de dise\~{n}o}

El erector es un sistema de pivote con un cilindro hidr\'{a}ulico de gas en la parte posterior del brazo. El cohete descansa sobre el brazo en orientaci\'{o}n horizontal, y el cilindro lo lleva a la vertical en un \'{u}nico movimiento continuo.

El pivote se ubica en el \textbf{centro de masa del veh\'{i}culo cargado} ($\approx 2500\ \text{mm}$ desde la tobera). Esta decisi\'{o}n minimiza el momento gravitacional que debe vencer el cilindro, reduciendo la fuerza requerida y el tama\~{n}o del sistema hidr\'{a}ulico.

\subsection{An\'{a}lisis de fuerzas}

Con el pivote en el CG, el momento gravitacional neto es idealmente cero. En la pr\'{a}ctica, la incertidumbre en la posici\'{o}n del CG y la masa del propio brazo del erector generan un momento residual. Estimando un desfase m\'{a}ximo de $\Delta_{CG} = 200\ \text{mm}$ entre el pivote y el CG real:

\begin{equation}
  M_{residual} = m_0 \cdot g \cdot \Delta_{CG}
  = 369 \times 9.81 \times 0.20 = \mathbf{724\ \text{N·m}}
\end{equation}

El cilindro hidr\'{a}ulico opera con el brazo de momento $d_{cil}$ desde el pivote. Con el cilindro ubicado a $d_{cil} = 1500\ \text{mm}$ del pivote (extremo posterior del brazo):

\begin{equation}
  F_{cil,min} = \frac{M_{residual}}{d_{cil} \cdot \sin\theta_{max}}
  = \frac{724}{1.5 \times 1.0} = \mathbf{483\ \text{N}}
\end{equation}

Aplicando un factor de seguridad de 3.0 (considera inercia rotacional, fricci\'{o}n en el pivote y viento):

\begin{equation}
  F_{cil,dise\~{n}o} = 483 \times 3.0 = \mathbf{1{,}450\ \text{N}}
\end{equation}

Un cilindro de gas est\'{a}ndar de \textbf{40 mm de di\'{a}metro a 50 bar} entrega:

\begin{equation}
  F_{cil,real} = P \cdot \frac{\pi d^2}{4}
  = 50 \times 10^5 \times \frac{\pi \times 0.04^2}{4}
  = \mathbf{6{,}283\ \text{N}}
\end{equation}

Factor de seguridad efectivo: $6{,}283 / 483 = \mathbf{13.0}$, que provee margen m\'{a}s que suficiente para un cilindro de catálogo est\'{a}ndar.

\subsection{Desplazamiento del CG durante el consumo de propelente}

El CG del veh\'{i}culo se desplaza durante la misi\'{o}n a medida que se consume el propelente. Este dato es relevante para la estabilidad en la plataforma y para el dise\~{n}o del brazo del erector:

\begin{equation}
  \Delta_{CG} = \frac{m_{prop} \times L_{motor}/2}{m_0}
  = \frac{119 \times 1{,}750}{369} \approx \mathbf{565\ \text{mm}}
\end{equation}

El CG se desplaza 565 mm hacia la nariz entre el estado lleno (despegue) y el estado vac\'{i}o (burnout). El erector opera con el veh\'{i}culo en estado lleno, por lo que el pivote debe dise\~{n}arse para la posici\'{o}n del CG lleno.

\subsection{Especificaciones del erector}

\begin{longtable}{L{5.5cm} C{3cm} L{5cm}}
  \rowcolor{tlvblue}
  \color{white}\textbf{Par\'{a}metro} &
  \color{white}\textbf{Valor} &
  \color{white}\textbf{Notas} \\
  \endfirsthead
  \rowcolor{tlvblue}
  \color{white}\textbf{Par\'{a}metro} &
  \color{white}\textbf{Valor} &
  \color{white}\textbf{Notas} \\
  \endhead
  \rowcolor{tlvgray}
  Tipo                        & Pivote + cilindro de gas  & Simple efecto \\
  Posici\'{o}n del pivote     & CG lleno ($\approx$2500 mm desde tobera) & Minimiza fuerza requerida \\
  \rowcolor{tlvgray}
  Brazo del erector           & 5200 mm                   & Longitud cohete + margen de apoyo \\
  Di\'{a}metro cilindro       & 40 mm                     & Componente de cat\'{a}logo \\
  \rowcolor{tlvgray}
  Presi\'{o}n de trabajo      & 50 bar                    & Est\'{a}ndar industrial \\
  Fuerza del cilindro         & 6{,}283 N                 & Factor seguridad efectivo 13.0 \\
  \rowcolor{tlvgray}
  Posici\'{o}n cilindro       & 1500 mm post-pivote       & Extremo posterior del brazo \\
  Material del brazo          & Perfil HEB-100 acero S275 & Rigidez bajo carga distribuida \\
  \rowcolor{tlvgray}
  Velocidad de ereci\'{o}n    & Controlada, $\leq$5°/s    & V\'{a}lvula de estrangulamiento \\
  Tiempo de ereci\'{o}n       & $\approx$18--20 s         & $90° / 5°/s$ \\
\end{longtable}

\subsection{Esquema del erector}

\begin{center}
\begin{tikzpicture}[scale=0.85]

  % Suelo
  \fill[gray!20] (-1, -0.3) rectangle (10, 0);
  \draw[gray, thick] (-1, 0) -- (10, 0);
  \node[font=\tiny, color=gray] at (9.5, -0.18) {Suelo};

  % Plataforma base
  \fill[steelgray!40] (3.5, 0) rectangle (6.5, 0.4);
  \draw[steelgray, thick] (3.5, 0) rectangle (6.5, 0.4);

  % Brazo del erector (posición inclinada ~45°, durante erección)
  \draw[tlvblue, line width=4pt]
    (5.0, 0.4) -- ++(45:5.3);
  % Pivote
  \fill[tlvred] (5.0, 0.4) circle (0.15);
  \node[font=\tiny, color=tlvred] at (4.4, 0.55) {Pivote (CG)};

  % Cohete sobre el brazo
  \draw[tlvmid, line width=10pt, opacity=0.4]
    (5.0, 0.4) -- ++(45:5.0);
  \draw[tlvmid, line width=2pt]
    (5.0, 0.4) -- ++(45:5.0);

  % Nariz del cohete
  \fill[tlvblue]
    ($(5.0,0.4)+(45:5.0)$) circle (0.12);

  % Cilindro hidráulico (posterior al pivote, ángulo ~20° desde horizontal)
  \draw[tlvorange, line width=3pt]
    (5.0, 0.4) -- ++(225:1.5)
    node[midway, below, font=\tiny, color=tlvorange] {};
  \draw[tlvorange, line width=3pt, dashed]
    ($(5.0,0.4)+(225:1.5)$) -- ++(270:1.2);
  \fill[tlvorange] ($(5.0,0.4)+(225:1.5)$) circle (0.12);
  \node[font=\tiny, color=tlvorange, align=left] at (2.5, -0.8)
    {Cilindro gas\\40mm / 50 bar};

  % Estado vertical (fantasma)
  \draw[tlvmid, line width=2pt, dashed, opacity=0.4]
    (5.0, 0.4) -- (5.0, 5.4);
  \draw[tlvblue, dashed, opacity=0.4]
    (5.0, 5.4) -- (5.0, 5.9);

  % Arco de erección
  \draw[gray, thin, -latex]
    ($(5.0,0.4)+(45:2.5)$) arc (45:88:2.5);
  \node[font=\tiny, color=gray] at (6.8, 2.8) {90°};

  % Cotas
  \draw[latex-latex, gray, thin]
    (5.0, -0.6) -- ++(45:1.5);
  \node[font=\tiny, color=gray] at (6.0, -0.3) {$d_{cil}=1500$ mm};

\end{tikzpicture}
\end{center}

% ── 5. ANÁLISIS DE RIESGO CINÉTICO Y ZONA NOTAM ──────────────────────────────
\section{An\'{a}lisis de Riesgo Cin\'{e}tico y Zona NOTAM}

\subsection{Escenario de peor caso}

El escenario que define el radio m\'{a}ximo de la zona de exclusi\'{o}n es el \textbf{fallo total en el instante de burnout}: FTS no activa, control de vuelo perdido, cohete a $v = 364.4\ \text{m/s}$ a $h = 6{,}248\ \text{m}$ con trayectoria horizontal.

\subsection{Energ\'{i}a cin\'{e}tica m\'{a}xima}

\begin{equation}
  E_{k,max} = \frac{1}{2} m_{dry} \cdot v_{burnout}^2
  = \frac{1}{2} \times 250 \times 364.4^2
  = \mathbf{16{,}597\ \text{kJ} \approx 16.6\ \text{MJ}}
\end{equation}

Equivalente en TNT: $16.6 / 4.184 \approx \mathbf{3.97\ \text{kg TNT}}$

\subsection{Alcance m\'{a}ximo en trayectoria balística}

Tiempo de ca\'{i}da libre desde $h_{burnout} = 6{,}248\ \text{m}$:

\begin{equation}
  t_{impacto} = \sqrt{\frac{2 \cdot h_{burnout}}{g}} = \sqrt{\frac{2 \times 6{,}248}{9.81}} = \mathbf{35.7\ \text{s}}
\end{equation}

Alcance horizontal sin arrastre (cota superior conservadora):

\begin{equation}
  x_{max} = v_{burnout} \cdot t_{impacto} = 364.4 \times 35.7 = \mathbf{13{,}005\ \text{m}}
\end{equation}

Alcance con arrastre aerodin\'{a}mico ($C_D = 0.35$, factor 0.65):

\begin{equation}
  x_{drag} = 0.65 \times 13{,}005 = \mathbf{8{,}453\ \text{m}}
\end{equation}

Alcance con empuje residual del motor (factor 1.20 sobre $x_{drag}$):

\begin{equation}
  x_{thrust} = 1.20 \times 8{,}453 = \mathbf{10{,}144\ \text{m}}
\end{equation}

\subsection{Radio NOTAM adoptado}

\begin{longtable}{L{6cm} C{2.5cm} C{2.5cm}}
  \rowcolor{tlvblue}
  \color{white}\textbf{Escenario} &
  \color{white}\textbf{Alcance (m)} &
  \color{white}\textbf{Cubierto por NOTAM} \\
  \endfirsthead
  \rowcolor{tlvblue}
  \color{white}\textbf{Escenario} &
  \color{white}\textbf{Alcance (m)} &
  \color{white}\textbf{Cubierto por NOTAM} \\
  \endhead
  \rowcolor{tlvgray}
  Trayectoria balística sin arrastre & 13{,}005 & \cmark \\
  Trayectoria balística con arrastre & 8{,}453  & \cmark \\
  \rowcolor{tlvgray}
  Con empuje residual (peor caso)    & 10{,}144 & \cmark \\
  \rowcolor{tlvlight}
  \textbf{Radio NOTAM adoptado}      & \textbf{11{,}000} & \textbf{\cmark\ con margen} \\
\end{longtable}

El radio de \textbf{11 km} cubre todos los escenarios con un margen de $\approx$8.4\% sobre el peor caso absoluto (trayectoria sin arrastre). La zona es una circunferencia conc\'{e}ntrica en la plataforma de lanzamiento.

\begin{center}
\begin{tikzpicture}[scale=0.42]
  \def\Rnotam{11}
  \def\Rthrust{10.144}
  \def\Rdrag{8.453}
  \def\Rnoatm{13.005}

  % Fondo
  \fill[tlvgray!30] (0,0) circle (\Rnotam);

  % Círculos
  \draw[gray!50, dashed, thick] (0,0) circle (\Rnoatm);
  \draw[tlvorange, thick] (0,0) circle (\Rdrag);
  \draw[tlvred, thick] (0,0) circle (\Rthrust);
  \draw[tlvblue, very thick] (0,0) circle (\Rnotam);

  % Plataforma
  \fill[tlvred] (0,0) circle (0.25);
  \node[font=\small\bfseries, color=tlvred] at (0, -0.7) {Plataforma};

  % Etiquetas
  \node[font=\small, color=gray!70] at (0, \Rnoatm+0.6)
    {Sin arrastre 13.0 km};
  \node[font=\small, color=tlvorange] at (0, \Rdrag+0.6)
    {Con arrastre 8.5 km};
  \node[font=\small, color=tlvred] at (0, -\Rthrust-0.6)
    {Empuje residual 10.1 km};
  \node[font=\small\bfseries, color=tlvblue] at (0, -\Rnotam-0.7)
    {NOTAM 11 km};

  % Cotas
  \draw[latex-latex, gray] (0,0) -- (\Rnotam, 0);
  \node[font=\small, color=gray] at (\Rnotam/2, 0.5) {11 km};

\end{tikzpicture}
\end{center}

% ── 6. SECUENCIA DE OPERACIONES DE LANZAMIENTO ───────────────────────────────
\section{Secuencia de Operaciones de Lanzamiento}

\begin{longtable}{C{1.2cm} L{4cm} C{2cm} L{6.5cm}}
  \rowcolor{tlvblue}
  \color{white}\textbf{Paso} &
  \color{white}\textbf{Actividad} &
  \color{white}\textbf{Tiempo} &
  \color{white}\textbf{Notas} \\
  \endfirsthead
  \rowcolor{tlvblue}
  \color{white}\textbf{Paso} &
  \color{white}\textbf{Actividad} &
  \color{white}\textbf{Tiempo} &
  \color{white}\textbf{Notas} \\
  \endhead
  \rowcolor{tlvgray}
  L-01 & Traslado AIT $\rightarrow$ plataforma & T$-$4 h & Veh\'{i}culo horizontal sobre carro de transporte \\
  L-02 & Acople al brazo del erector & T$-$3.5 h & Fijaci\'{o}n mec\'{a}nica al brazo en posici\'{o}n horizontal \\
  \rowcolor{tlvgray}
  L-03 & Ereci\'{o}n a la vertical & T$-$3 h & Cilindro hidr\'{a}ulico; velocidad $\leq 5°/s$ \\
  L-04 & Fijaci\'{o}n a la plataforma & T$-$2.5 h & Pernos de retenci\'{o}n al riel de la plataforma \\
  \rowcolor{tlvgray}
  L-05 & Conexi\'{o}n umbilicales & T$-$2 h & Alimentaci\'{o}n el\'{e}ctrica y telemetr\'{i}a de tierra \\
  L-06 & Verificaci\'{o}n final de sistemas & T$-$1 h & Check aviónica, FTS (pasador a\'{u}n insertado), telemetr\'{i}a \\
  \rowcolor{tlvgray}
  L-07 & Despeje de personal & T$-$15 min & Todo el personal fuera del radio de seguridad inmediato \\
  L-08 & Retiro de pasador FTS & T$-$5 min & \'{U}nico operador autorizado; retiro y confirmaci\'{o}n verbal \\
  \rowcolor{tlvgray}
  L-09 & Desconexión umbilicales & T$-$10 s & Autom\'{a}tica por comando del AFCC \\
  L-10 & Ignición & T$-$0 & Secuencia: pirot\'{e}cnico $\rightarrow$ igniter s\'{o}lido $\rightarrow$ motor principal \\
\end{longtable}

% ── 7. ELEMENTOS PENDIENTES ──────────────────────────────────────────────────
\section{Elementos Pendientes}

\begin{longtable}{C{1cm} L{8cm} C{2.2cm}}
  \rowcolor{tlvblue}
  \color{white}\textbf{ID} &
  \color{white}\textbf{Descripci\'{o}n} &
  \color{white}\textbf{Prioridad} \\
  \endfirsthead
  \rowcolor{tlvblue}
  \color{white}\textbf{ID} &
  \color{white}\textbf{Descripci\'{o}n} &
  \color{white}\textbf{Prioridad} \\
  \endhead
  \rowcolor{tlvgray}
  P-01 & Determinar posici\'{o}n exacta del CG lleno mediante modelado CAD; ajustar posici\'{o}n del pivote. & \textcolor{tlvred}{\textbf{Alta}} \\
  P-02 & Dise\~{n}ar sistema de retenci\'{o}n del cohete en el brazo del erector durante el traslado y la ereci\'{o}n. & \textcolor{tlvred}{\textbf{Alta}} \\
  \rowcolor{tlvgray}
  P-03 & Dise\~{n}ar sistema de desconexión autom\'{a}tica de umbilicales en T$-$10 s. & \textcolor{tlvorange}{\textbf{Media}} \\
  P-04 & Verificar normativa local para zona NOTAM de 11 km (autoridad de aviaci\'{o}n civil competente). & \textcolor{tlvorange}{\textbf{Media}} \\
  \rowcolor{tlvgray}
  P-05 & Especificar carro de transporte AIT $\rightarrow$ plataforma compatible con brazo del erector. & \textcolor{tlvorange}{\textbf{Media}} \\
  P-06 & Validar posici\'{o}n del CG vac\'{i}o post-burnout para estabilidad en trayectoria. & \textbf{Baja} \\
\end{longtable}

\end{document}